\chapter{弃牌与跟注}

\section{本章导入}

本章讨论弃牌与跟注。过度跟注是娱乐玩家最常见、也最昂贵的错误之一。跟注不应是好奇、侥幸或不甘心的情绪行为，而应当是基于底池赔率、所需胜率和对手范围的数学行为。

学会弃牌并不代表软弱，而是尊重价格和风险。稳定的弃牌纪律可以让新手避免在明显劣势局面中持续扩大损失。

\section{核心概念}

本章将介绍底池赔率、所需胜率、对手范围、抓诈唬、阻断牌和弃牌纪律。读者需要学会把“我不想被诈唬”转化为“这个价格下，我的牌是否有足够胜率”。

\section{新手常见误区}

常见误区包括不计算赔率、面对大注用情绪跟注、只想着对手可能在诈唬、忽略对手价值组合数量，以及把弃牌视为失败。

\section{实战思考框架}

面对下注时，先计算跟注所需胜率，再估算对手价值牌和诈唬牌比例。如果对手的诈唬不足以支撑价格，正确选择就是弃牌。

\section{牌例拆解}

\subsection{牌例一：待补充}

本牌例将用于计算面对半池下注时的跟注所需胜率。

\subsection{牌例二：待补充}

本牌例将用于分析河牌抓诈唬时需要满足的条件。

\section{本章总结}

弃牌与跟注是风险管理的核心。新手应把跟注从情绪行为改造成数学行为，并通过弃牌纪律保护自己的长期 EV。

\section{课后训练}

请记录最近五次面对大注的决策，分别写出底池大小、下注大小、跟注所需胜率和最终行动。
